% 方法章节草稿：根据用户于 2026-10-19 提供的文本保存。
% 该文件为独立草稿，不包含导言区；可由主 TeX 文件通过 \input{} 引入。
\section{方法}
\subsection{系统架构：事件--证据--行动的异常预检管线}

本研究使用事件--证据--行动管线，将原始中子电流 CSV 数据转化为可审查的人工核查建议。如图 X 所示，该管线包含三个阶段：（i）时序事件解析与事件卡构建；（ii）人工证据卡组织与候选证据获取；（iii）证据条件下的受限行动选择。在阶段 1，原始中子电流记录被确定性地解析为原子事件、单通道事件卡和聚合事件组合。在阶段 2，既定人工证据卡库保留证据文本、来源、章节或行号、主题标签及统一适用性标签；系统根据实验配置使用 BM25 获取候选证据，或由 LLM 面向全库选择候选证据。在阶段 3，事件组合、候选证据与固定检查建议定义共同构成推理上下文，语言模型据此输出一项人工检查建议。由于最终建议质量共同取决于事件表示、证据获取和行动选择，本文以第 3.2 节定义的指标评价端到端流程，并在实验部分报告结果。

\textbf{阶段 1：时序事件解析与事件卡构建。}
令
\[
\mathcal{S}=\{s_1,s_2,\ldots,s_N\}
\tag{1}
\]
表示输入的中子电流 CSV 记录集合，其中每条记录包含时间、通道和对应的电流观测值。系统对每个通道的时序记录采用确定性分析规则，识别具有预定义模式的局部异常，并将其表示为原子事件：
\[
\mathcal{U}=\{u_i=f_{\mathrm{event}}(s_i)\mid s_i\in\mathcal{S}\}.
\tag{1}
\]
原子事件保留其发生时间、所属通道、异常模式及相关数值摘要。随后，系统将同一通道内具有时间关联的原子事件组织为单通道事件卡：
\[
c^{\mathrm{single}}_j=f_{\mathrm{single}}(\mathcal{U}_j),
\tag{2}
\]
并进一步将跨通道或相邻时间段内相关的单通道事件卡组合为聚合事件组合：
\[
x_l=f_{\mathrm{group}}\bigl(\{c^{\mathrm{single}}_j\}\bigr).
\tag{3}
\]
原子事件、单通道事件卡与聚合事件组合均作为独立但相互关联的对象保存。借助共享的时间范围、通道标识和事件关系，该表示既保留原始异常的可回溯性，也为后续证据获取提供统一的文本化任务输入。

\textbf{阶段 2：人工证据卡组织与候选证据获取。}
令
\[
\mathcal{Y}=\{y_1,y_2,\ldots,y_M\}
\tag{}
\]
表示人工整理的证据卡库；本研究中 \(M=300\)。每张证据卡 \(y_j\) 包含证据文本、来源信息、章节或行号、主题标签和统一适用性标签。与自动文档分块不同，证据卡是既定的人工证据单元，不在实验过程中按段落重新切分。对于采用 BM25 的 B0 和 B1 配置，系统将事件组合 \(x\) 转换为检索查询 \(q(x)\)，并计算其与各证据卡之间的词汇相关性：
\[
r_j=\mathrm{BM25}\bigl(q(x),y_j\bigr).
\tag{4}
\]
检索器据此返回得分最高的 Top-K 候选证据集合：
\[
\mathcal{Y}_K(x)=\mathrm{TopK}_{y_j\in\mathcal{Y}}\, r_j.
\tag{5}
\]
在 B0 中，该候选集合进一步依据证据卡的适用性标签进行机械投票；在 B1 中，同一候选集合被提供给 LLM 进行行动选择。对于 B2，系统不使用预先固定的 BM25 候选包，而是向 LLM 提供完整证据库 \(\mathcal{Y}\)，由其针对当前事件组合独立选择指定数量的候选证据：
\[
\mathcal{Y}_K(x)=
f_{\mathrm{select}}\bigl(x,\mathcal{Y},K\bigr).
\tag{6}
\]
因此，B1 用于考察固定检索候选条件下的证据--行动推理，B2 则同时考察全库证据选择与行动选择的联合表现。

\textbf{阶段 3：证据条件下的受限行动选择。}
候选证据集合 \(\mathcal{Y}_K(x)\) 与事件组合 \(x\)、固定行动空间 \(\mathcal{A}\) 以及系统约束提示 \(P_{\mathrm{sys}}\) 共同构成行动选择上下文，其中
\[
\mathcal{A}=
\{\mathrm{A1},\mathrm{A2},\mathrm{A3},\mathrm{A4}\}.
\]
四项行动分别对应现场操作或校准核查、跨通道相似异常比较、记录处理与时间戳核查，以及信号传输路径核查。语言模型在可选的固定示例上下文 \(P_{\mathrm{ex}}\) 条件下输出一项检查建议：
\[
\hat{z}=
f_{\mathrm{LLM}}
\bigl(x,\mathcal{Y}_K(x),\mathcal{A},P_{\mathrm{sys}},P_{\mathrm{ex}}\bigr),
\quad \hat{z}\in\mathcal{A}.
\tag{7}
\]
在 B2 中，证据选择和行动选择是针对每个案例独立执行的两个阶段；行动选择阶段只使用该次证据选择所产生的候选证据，不复用其他案例、其他配置或其他运行中的选择结果。输出以 JSON 形式返回，并接受格式与内容合同验证：
\[
V(\hat{z})=
\begin{cases}
1, & \text{若输出合法且仅包含一项 } \mathcal{A} \text{ 中的行动；}\\
0, & \text{否则。}
\end{cases}
\tag{8}
\]
\subsection{定量评价}

现有检索增强与异常分析评价通常强调候选文档排序、字符串匹配或最终分类准确率，却较少同时检验候选证据是否与人工检查意图一致、证据是否足以支持相应的核查建议，以及系统输出是否遵守预先定义的行动边界。常用的 Precision@K、Recall@K 和分类 Accuracy 等指标对于通用检索或分类任务具有参考价值，却不能直接评价不同证据卡在语义上是否承担相同的核查支持作用，也不能区分“未获取合适证据”与“已获取合适证据但行动选择错误”这两类失效模式。对于中子电流异常预检而言，这种区分尤为重要：系统目标不是确认具体设备故障，而是在可追溯证据支持下，为工程人员提出一项固定的优先核查建议。

因此，本文引入一组面向事件--证据--行动工作流的评价指标，用于刻画候选证据质量、证据与行动之间的一致性、最终行动选择正确性以及输出格式合规性。正式评价集由冻结的 \(N=60\) 个金标准案例构成：
\[
\mathcal{G}_{\mathrm{eval}}=
\{(x_i,\mathcal{Y}^{*}_i,\mathcal{U}^{*}_i,a^{*}_i)\}_{i=1}^{N}.
\tag{9}
\]
其中，\(x_i\) 为事件组合，\(\mathcal{Y}^{*}_i\) 为人工整理的金标准证据卡集合，\(\mathcal{U}^{*}_i\) 为该证据集合对应的金标准语义单元集合，\(a^{*}_i\) 为固定人工检查建议。每个金标准案例包含一个由 1--4 个事件原子构成的事件组合、3 张具有同一检查意图的证据卡，以及 A1--A4 中的一项固定检查建议。金标准证据集合并不假定穷尽证据库中所有可能有效的支持来源；它提供评价与人工核验的统一锚点。对每个案例，系统产生候选证据集合 \(\hat{\mathcal{Y}}_{i,K}\) 及行动预测 \(\hat{a}_i\)。以下指标逐案例计算后，再在冻结评估集上汇总。金标准案例标识仅用于离线评价，且不传入任何模型提示词。

\textbf{1. 语义证据精确度、召回率与 F1 值。}
候选证据质量不能仅依据证据卡 ID 是否与金标准卡完全一致来判断，因为不同证据卡可能表达相同或高度相近的核查支持语义。令 \(m(\mathcal{Y})\) 表示由既定语义单元映射获得的证据语义集合，候选集合的语义表示为 \(\hat{\mathcal{U}}_{i,K}=m(\hat{\mathcal{Y}}_{i,K})\)。语义证据精确度（Semantic Precision@K，SP@K）衡量前 \(K\) 张候选证据中，与金标准证据语义一致的部分所占比例：
\[
\mathrm{SP@K}_i=
\frac{|\hat{\mathcal{U}}_{i,K}\cap\mathcal{U}^{*}_i|}
{\max(|\hat{\mathcal{U}}_{i,K}|,1)}.
\tag{10}
\]
语义证据召回率（Semantic Recall@K，SR@K）衡量金标准所要求的证据语义中，被前 \(K\) 张候选证据覆盖的比例：
\[
\mathrm{SR@K}_i=
\frac{|\hat{\mathcal{U}}_{i,K}\cap\mathcal{U}^{*}_i|}
{|\mathcal{U}^{*}_i|}.
\tag{11}
\]
语义证据 F1 值（Semantic F1@K，SF1@K）为精确度与召回率的调和平均：
\[
\mathrm{SF1@K}_i=
\frac{2\cdot\mathrm{SP@K}_i\cdot\mathrm{SR@K}_i}
{\mathrm{SP@K}_i+\mathrm{SR@K}_i},
\tag{12}
\]
当分母为零时定义为零。较高的 SP@K 表示候选包较少包含与当前检查意图无关的证据；较高的 SR@K 表示候选包能够覆盖更多人工金标准所要求的证据语义；SF1@K 则综合衡量证据选择的准确性与完整性。

\textbf{2. 语义候选覆盖率。}
语义候选覆盖率（Semantic Candidate Coverage@K，SCC@K）衡量每个案例的候选证据包是否至少包含一项与金标准一致的证据语义：
\[
\mathrm{SCC@K}_i=
\begin{cases}
1, & \text{若 } \hat{\mathcal{U}}_{i,K}\cap\mathcal{U}^{*}_i\neq\varnothing;\\
0, & \text{否则。}
\end{cases}
\tag{13}
\]
该指标不同于 SR@K：SR@K 关注金标准语义被覆盖的程度，而 SCC@K 仅回答候选包是否为后续行动选择提供了至少一项可接受的证据依据。其系统级均值表示在全部案例中，系统成功向后续阶段提供可接受候选证据的比例。

\textbf{3. 语义--行动对 F1 值。}
仅有语义正确的候选证据并不保证最终建议正确；同样，偶然预测正确的行动也不代表模型建立了正确的证据--行动关联。为同时评价两者，本文构造候选证据语义与行动标签组成的集合：
\[
\hat{\mathcal{P}}_{i,K}=
\{(u,\hat{a}_i)\mid u\in\hat{\mathcal{U}}_{i,K}\},
\quad
\mathcal{P}^{*}_i=
\{(u,a^{*}_i)\mid u\in\mathcal{U}^{*}_i\}.
\]
语义--行动对精确度和召回率分别定义为候选对集合与金标准对集合的交集比例，并据此计算语义--行动对 F1 值（Semantic Candidate--Action Pair F1@K，SCA-F1@K）：
\[
\mathrm{SCA\mbox{-}F1@K}_i=
\frac{2P^{\mathrm{pair}}_{i,K}R^{\mathrm{pair}}_{i,K}}
{P^{\mathrm{pair}}_{i,K}+R^{\mathrm{pair}}_{i,K}}.
\tag{14}
\]
该指标要求候选证据语义与最终检查建议同时正确；若候选证据语义正确但行动错误，或行动正确但候选证据未覆盖金标准语义，则相应证据--行动对均不计为命中。

\textbf{4. 最终行动准确率与宏平均 F1 值。}
最终行动准确率（Action Accuracy，Acc）衡量系统是否预测出与金标准一致的固定人工检查建议：
\[
\mathrm{Acc}=
\frac{1}{N}\sum_{i=1}^{N}
\mathbb{I}(\hat{a}_i=a^{*}_i).
\tag{15}
\]
由于四类检查建议在评估集中可能并非完全均衡，本文同时报告宏平均 F1 值（Macro-F1，MF1）：
\[
\mathrm{MF1}=
\frac{1}{|\mathcal{A}|}
\sum_{a\in\mathcal{A}}\mathrm{F1}(a),
\quad
\mathcal{A}=\{\mathrm{A1},\mathrm{A2},\mathrm{A3},\mathrm{A4}\}.
\tag{16}
\]
其中，\(\mathrm{F1}(a)\) 为将行动 \(a\) 视为正类时得到的类别 F1 值。Macro-F1 对四类行动给予相同权重，因此能够避免高频行动类别主导总体结果。

\textbf{5. JSON 合法率与输出合同合规率。}
除证据与行动质量外，系统还必须满足可机器读取和任务边界要求。JSON 合法率（JSON Validity Rate，JVR）衡量模型输出能否被正确解析为预期 JSON 结构：
\[
\mathrm{JVR}=
\frac{1}{N}\sum_{i=1}^{N}
\mathbb{I}(\mathrm{parse}(\hat{o}_i)=1),
\tag{17}
\]
其中 \(\hat{o}_i\) 为案例 \(i\) 的原始模型输出。输出合同合规率（Output Contract Compliance Rate，OCCR）进一步衡量解析后的输出是否满足当前阶段的全部约束：
\[
\mathrm{OCCR}=
\frac{1}{N}\sum_{i=1}^{N}
\mathbb{I}(\hat{o}_i\text{ satisfies contract}).
\tag{18}
\]
对于行动选择阶段，合同要求输出恰好一项 A1--A4 中的检查建议；对于 B2 证据选择阶段，合同还要求输出恰好 \(K\) 张来自既定证据库、且互不重复的候选证据。JSON 合法率只检验格式可解析性，而合同合规率检验输出是否同时满足结构、字段和值域等任务约束。

\textbf{平均指标。}
系统级表现报告为冻结 60 例评估集上各逐案例指标的均值：
\[
\overline{M}=
\frac{1}{N}\sum_{i=1}^{N}M_i.
\tag{19}
\]
其中，SCC@K 的均值表示至少获得一项可接受候选证据的案例比例；Acc 表示最终检查建议正确的案例比例；JVR 和 OCCR 分别表示输出在语法层面与任务合同层面的合规比例。这些指标为比较 BM25 与适用性标签机械投票基线、固定 BM25 候选包上的 LLM 行动选择，以及全证据库 LLM 证据选择与行动选择提供了紧凑框架。它们不将系统表现压缩为单一分数，而是分别衡量：（i）候选证据的语义准确性与完整性；（ii）证据语义与检查建议的联合一致性；（iii）最终行动选择的正确性；以及（iv）输出格式和任务边界的合规性。

\textbf{语义单元映射与严格证据卡 ID 匹配。}
在指标设计中，本文不采用“候选证据卡 ID 必须与金标准证据卡 ID 完全一致”的严格逐卡匹配作为唯一评价依据。实践中，同一检查意图可能由多张来自不同来源、但表达相近技术依据的证据卡支持；若只按卡片 ID 评分，语义上可接受的替代证据会被错误地判为未命中。相反，本文使用人工建立的证据卡--语义单元映射，将候选证据与金标准证据在检查支持语义层面进行比较。该映射并不把所有相似文本视为等价，而是仅将能够支持同一人工核查意图的既定证据单元归入相应语义类别。由此，SF1@K、SCC@K 和 SCA-F1@K 能够在保留金标准约束的同时，避免评价结果被表面卡片标识或措辞差异主导。

\subsection{系统模块}

事件--证据--行动管线以原始中子电流 CSV 的确定性事件解析和人工证据卡库为主要输入机制，并通过表~\ref{tab:modules} 概括的功能模块完成候选证据获取、受限行动选择和结果验证。与由代理根据开放式用户查询在运行时任意选择工具的系统不同，本文的模块调用顺序由实验配置预先确定：B0 使用 BM25 检索与适用性标签机械投票；B1 使用 BM25 候选证据包与 LLM 行动选择；B2 则依次执行全库 LLM 证据选择与 LLM 行动选择。该设计使不同配置能够在同一事件表示、同一人工证据库和同一冻结评估集上进行可复现比较，同时保持候选证据、行动输出与评价结果之间的可追溯关系。

\begin{table}[t]
\centering
\small
\caption{中子电流异常预检管线的功能模块。}
\label{tab:modules}
\setlength{\tabcolsep}{3pt}
\begin{tabular}{p{0.30\columnwidth} p{0.62\columnwidth}}
\hline
\textbf{模块} & \textbf{功能} \\
\hline
CSV 事件解析模块 & 对原始中子电流 CSV 执行确定性分析，识别原子事件。 \\
事件卡构建模块 & 将原子事件组织为单通道事件卡和聚合事件组合。 \\
人工证据卡库 & 保存来自 39 份资料来源的 300 张人工证据卡及其来源、定位信息、主题标签和适用性标签。 \\
BM25 候选证据获取模块 & 按词汇相关性返回 Top-K 候选证据，用于 B0 和 B1。 \\
标签投票模块 & 依据候选证据的标签对 A1--A4 进行机械投票，构成 B0 基线。 \\
LLM 证据选择模块 & 从完整证据库中独立选择 \(K\) 张候选证据，用于 B2。 \\
LLM 行动选择模块 & 根据事件组合和候选证据，在 A1--A4 中选择一项人工检查建议。 \\
输出与统计模块 & 检查 JSON 与输出合同，并汇总证据质量和行动选择指标。 \\
\hline
\end{tabular}
\end{table}

该框架独立于某一特定预训练模型的内部参数：模型在所有实验中保持冻结，使用远程 LLM API 时，系统会将当前实验配置所需的事件组合、候选证据或全库证据卡及提示内容提供给模型。

在方法执行层面，原始 CSV、事件卡、证据卡、金标准案例、个案 JSON/JSONL 输出及汇总量化指标均保持为彼此可关联的研究对象。对于 LLM 输出，系统分别检查 JSON 是否可解析，以及输出是否满足当前阶段的要求，例如行动选择是否只返回一项 A1--A4 行动，或证据选择是否返回指定数量且来自既定证据库的候选卡片。通过来源定位、个案记录和合同检查，系统将异常事件、证据依据和人工检查建议之间的关系保留为可逐项审查的链条。

\subsection{实验配置、敏感性分析与基准设计}

所有正式比较均使用冻结的 60 个案例，并依据第 3.2 节定义的金标准证据语义和固定行动标签评分。本文分别报告候选证据质量、语义--行动对质量、最终行动表现、JSON 合法率和输出合同合规率，而不将其合并为单一分数。

B0 使用 BM25 Top-K 检索与适用性标签机械投票；B1 在相同的 BM25 候选证据包上使用 LLM 选择行动；B2 则由 LLM 面向完整证据库选择候选证据，再据此选择行动。所有 LLM 实验使用参数冻结的 \texttt{qwen3.8-flash}，温度设为 0，并按案例串行执行。候选证据数量取 \(K\in\{3,5,8,10\}\)。B1 比较零样本与不同数量的固定示例，并对同一示例数量使用三个冻结示例池独立运行。

B2 分别在证据选择端（\textit{selector}）、行动选择端（\textit{action}）和两个阶段（\textit{both}）加入固定示例。每个案例均重新执行证据选择与行动选择，不复用其他运行的候选证据。实验保留个案 JSON/JSONL 输出、错误记录、token 使用量和延迟统计，并以 CSV 保存汇总结果；请求失败、解析失败或违反输出合同的结果不计为正常模型性能。
